\section{Trabajos relacionados}
\label{sec:trabajos_relacionados}

Este trabajo se sitúa en la intersección de cuatro cuerpos de literatura (calibración, reglas de puntuación
propias, predicción selectiva y calibración post-hoc) y de un quinto, reciente y aún en formación: el
ecosistema de modelos de decisión tipo Jev.

\textbf{Calibración de clasificadores.} Un predictor probabilístico está calibrado si, de las veces que
anuncia confianza $c$, acierta con frecuencia $c$. \cite{niculescu2005} mostró que los clasificadores
clásicos están sistemáticamente mal calibrados. \cite{guo2017} lo reabrió para las redes profundas modernas
y popularizó el \textbf{Expected Calibration Error (ECE)} con su \textbf{temperature scaling}. El ECE,
formalizado por \cite{naeini2015}, depende críticamente del esquema de partición (bins de igual ancho frente
a bins de igual masa), un detalle que el ecosistema Jev-OSS omite al reportar. La \textbf{discriminación}
(capacidad de ordenar aciertos por confianza, medida con AUROC) es una propiedad distinta de la calibración;
ambas se reportan aquí.

\textbf{Reglas de puntuación propias.} \cite{gneiting2007} caracterizan las reglas \textbf{estrictamente
propias}, cuyo óptimo se alcanza declarando la distribución verdadera. Es la base teórica de RLCD: la
recompensa de Jev combina la regla logarítmica con la esférica, y el RPS para salidas ordinales
\cite{good1952,epstein1969}. La teoría predice que reportar honestamente es la única estrategia óptima; la
práctica del ecosistema la contradice, como se ve en el cierre de esta sección.

\textbf{Predicción selectiva.} \cite{chow1970} derivó la política óptima de rechazo: actuar cuando
$p > 1 - C_{\mathrm{esc}}/C_{\mathrm{err}}$. \cite{geifman2017} llevaron la idea a redes profundas y a las
curvas \emph{risk--coverage} y el AURC; \cite{jitkrittum2023} estudian cuándo basta la confianza para
derivar. Todos estos umbrales suponen que la confianza está calibrada, que es exactamente lo que este
trabajo pone a prueba cuando la confianza tiene varias definiciones.

\textbf{Calibración post-hoc y conformal.} \cite{platt1999} y \cite{zadrozny2002} dan la escala de Platt y
la regresión isotónica; \cite{guo2017} la temperatura. \cite{vovk2005,angelopoulos2021} ofrecen, con la
predicción conformal, cobertura con garantía libre de distribución, como alternativa a calibrar
\texttt{confidence}.

\textbf{El ecosistema Jev-OSS (2026).} \cite{jev2026blog,jev2026docs} describen el modelo y sus primitivos;
\cite{archerhume,navinpai} reconstruyen su arquitectura y la mecánica de RLCD a partir de sondeos de la API.
\cite{layaRepo} publica la primera implementación abierta con pesos y receta;
\cite{kevRepo,nimble,semif,jeff} le siguen. Los cierres de septiembre de 2026 (por qué RLCD no calibra
\cite{arXiv2609.38850,lakoreai2026}, la coherencia de la negación \cite{arXiv2609.33209}, el sesgo
nombre--definición \cite{arXiv2609.26758} y la cartografía del ecosistema \cite{arXiv2609.30216}) ocupan el
\guillemotleft por qué\guillemotright{} y el \guillemotleft qué hay\guillemotright. La reproducción
pre-registrada de \cite{arXiv2609.33843} es el estado del arte más cercano: mide la calibración de Laya
\textbf{dentro de una sola definición} ($q = p_{\max}$) y encuentra que el umbral no transfiere \textbf{entre
datasets}. Lo que ninguno hace es comparar \textbf{entre definiciones}.

\textbf{El hueco.} La literatura mide la calibración de un modelo bajo una definición, o comprueba
cualitativamente que un umbral no viaja. \textbf{No existe la matriz que cuantifica, par a par y definición
a definición, cuánta cobertura y precisión se pierden al portar un umbral.} Esa es la contribución de este
trabajo.
